% Bewertungspaket zum Gesamttest «Lineare Gleichungssysteme» — Quelle des PDFs (python3 scripts/build-lp-pdf.py gleichungssysteme).
% Ganze Punkte je Teilschritt; Werte und Fehlerwerte mit python3 nachgerechnet (scripts/lp/gleichungssysteme/zahlen.py, 11.10.2026).
\documentclass[11pt]{article}
\usepackage{../lp-druck}
\renewcommand{\lpname}{Lineare Gleichungssysteme}
\renewcommand{\lpbereich}{Grundlagenfach}
\renewcommand{\lpteilgebiet}{2.3}
\renewcommand{\lpfuss}{Bewertungspaket}
\newcolumntype{P}{>{\centering\arraybackslash}p{10mm}}
\newenvironment{raster}{\par\smallskip\noindent\tabularx{\linewidth}{@{}>{\raggedright\arraybackslash}p{38mm}P X@{}}\toprule
  \small\textsc{Teilschritt} & \small\textsc{P} & \small\textsc{Musterlösung}\\\midrule}{\bottomrule\endtabularx\par}
\newcommand{\fehler}[1]{\par\smallskip{\small\textcolor{lprot}{\bfseries Typische Fehler:} #1}}
\newcommand{\E}{\,{\footnotesize\bfseries\color{lpgruen}(E)}}
\begin{document}
\lpkopf{Bewertungspaket zum Gesamttest}{}

\begin{lpachtung}{Erst lösen, dann öffnen}
Dieses Paket enthält die Musterlösung. Wer es vor dem Lösen liest, misst am Ende nur, wie gut das Abschreiben gelingt.
\end{lpachtung}

\section*{1 · So gehst du vor}
\begin{enumerate}[leftmargin=*,itemsep=2pt]
\item \textbf{Gesamttest lösen} — vollständig, mit Rechenweg, ohne dieses Paket und ohne Taschenrechner.
\item \textbf{Fotografieren oder scannen}, gut lesbar, eine Seite pro Bild. \textbf{Kein Name} auf den Bildern.
  Handyfotos tragen oft unsichtbar Standort, Zeit und Gerät mit (Metadaten). Lade darum lieber einen Scan oder ein
  Bildschirmfoto des Fotos hoch, oder entferne vorher die Standortangaben.
\item \textbf{Bewerten lassen.} Ob du dafür eine KI nutzt und welche, entscheidest du selbst und in eigener Verantwortung —
  je nachdem, was dir aufgrund deines Alters und deiner persönlichen Situation erlaubt ist (Altersgrenzen und
  Nutzungsbedingungen des Anbieters, allenfalls Einverständnis deiner Eltern). Lade nur deine Lösung und dieses PDF hoch,
  keine persönlichen Daten, und füge den Auftrag aus Abschnitt~2 ein. Ohne KI kannst du dich mit Abschnitt~3 selbst bewerten.
\item \textbf{Rückmeldung prüfen.} Stimmt, was die KI aus deiner Handschrift gelesen hat? Eine KI kann sich irren —
  im Zweifel gilt das Raster in Abschnitt~3.
\item \textbf{Gezielt wiederholen} nach Abschnitt~4.
\end{enumerate}

\section*{2 · Auftrag an die KI}
\begin{tcolorbox}[colback=lplinie!25,colframe=lplinie,boxrule=0.5pt,arc=1mm,fontupper=\small]
Du bist Korrektorin oder Korrektor für Mathematik an einer Schweizer Berufsmaturitätsschule. Im Anhang findest du (1)~das Bewertungspaket zum Gesamttest «Lineare Gleichungssysteme» mit Musterlösung und Punkteraster und (2)~Fotos meiner handschriftlichen Lösung.\par\medskip
Geh so vor:\par
1. Schreib zu jeder Aufgabe G1 bis G8 zuerst ab, was du in meiner Lösung liest — Rechenweg und Ergebnis. Was du nicht sicher lesen kannst, markierst du mit [unsicher]. Nicht raten. Bei G1 beschreibst du, durch welche Punkte meine Geraden gehen, soweit du es erkennst.\par
2. Vergib die Punkte genau nach dem Raster im Bewertungspaket (Abschnitt 3), ganze Punkte je Zeile. Ziehe einen Fehler nur einmal ab: Rechne ich mit einem falschen Zwischenergebnis richtig weiter, gibt es die Folgepunkte. Ausnahme: Zeilen mit \textbf{(E)} gibt es nur für das richtige Ergebnis, nie als Folgepunkt. Die Zeilen «Typische Fehler» sagen, welche Rasterzeile bei diesem Fehler 0 Punkte gibt — sie sind kein zusätzlicher Abzug.\par
3. Andere richtige Lösungswege zählen voll (ein anderes Verfahren, eine andere Variable zuerst eliminiert). Gleichwertige Schreibweisen zählen voll: Dezimalpunkt oder Dezimalkomma, Bruch oder Dezimalzahl ($\tfrac72 = 3.5$), $\mathbb{L} = \{(2 \mid 3)\}$ oder «$x = 2$, $y = 3$» oder $(2;\, 3)$, $\mathbb{L} = \{\,\}$ oder «keine Lösung» — ausser wo das Raster eine bestimmte Form verlangt. Zeilen mit \textbf{(E)} sind Ergebnispunkte für Ergebnisse direkt aus den Angaben: Das Ergebnis muss stimmen \emph{und} aus einem erkennbaren Weg hervorgehen. Zeilen mit «Folgezeile» bewerten Schritte aus eigenen Zwischenwerten: richtig weitergerechnet gibt den Punkt. Zeilen mit \textbf{(A)} sind Ablesepunkte: Der abgelesene Wert genügt. Alle übrigen Zeilen bewerten den Weg selbst.\par
4. Begründe jeden Abzug in einem Satz und nenne die Fehlerart (zum Beispiel «Klammer beim Einsetzen vergessen» oder «rechte Seite nicht vervielfacht»). Schreib die Musterlösung nicht ab — zeig mir, wo mein Fehler liegt.\par
5. Gib zum Schluss eine Tabelle «Aufgabe | Punkte | Begründung», die Gesamtpunktzahl (von 25) und — nach der Selbsteinschätzung in Abschnitt 4 — welches Kapitel des Leitprogramms ich wiederholen soll. Keine Note.\par
6. Fehlt ein Foto oder ist eine Aufgabe unlesbar, sag es, statt sie zu bewerten.
\end{tcolorbox}

\needspace{16\baselineskip}
\section*{3 · Musterlösung und Punkteraster}
\textbf{Allgemein:} Ganze Punkte je Zeile. Ein Fehler wird nur einmal abgezogen (Folgefehler): Zeilen ohne Zeichen geben den
Punkt auch, wenn mit einem eigenen falschen Zwischenwert richtig weitergerechnet wird. \textbf{(E)}~= Ergebnispunkt für ein Ergebnis direkt aus den Angaben, nur wenn es stimmt und der Weg erkennbar ist; \textbf{(A)}~= Ablesepunkt; «Folgezeile» = Schritt aus eigenen Zwischenwerten, richtig weitergerechnet gibt den Punkt. Derselbe Begriffsfehler in einer späteren Zeile derselben
Aufgabe gibt keinen zweiten Folgepunkt. Der ganze Test ist ohne Taschenrechner gestellt.
«Typische Fehler» nennt, welche Zeile 0 Punkte gibt, und die Punkte, die bleiben.
\textbf{Teilrichtige Zeilen:} Verlangt eine Zeile mehrere Angaben und stimmt nur ein Teil, gibt die Zeile 0 Punkte —
ausser die Zeile nennt ausdrücklich eine Aufteilung.
\textbf{Probe:} Probe heisst Einsetzen in die gegebenen Gleichungen (bei G5 auch die geordnete Form). Eine richtig gerechnete Probe zur eigenen falschen Lösung, die den Fehler aufdeckt (beide Seiten verschieden, so benannt), gibt den Probepunkt.

\begin{lpaufgabe}{G1}{Zeichnen und bestätigen}{3 P}
\begin{raster}
(a) Auflösen und zeichnen & 1 & (I) $y = -0.5x + 3$, zum Beispiel durch $(0 \mid 3)$ und $(6 \mid 0)$; (II) $y = x - 3$, durch $(0 \mid -3)$ und $(3 \mid 0)$. Beide Geraden nötig. Wer beim Auflösen einen Fehler macht und die eigene Gerade richtig zeichnet: 0 (die Zeile verlangt beides).\\
(b) Ablesen & 1\,{\footnotesize\bfseries\color{lpgruen}(A)} & Schnittpunkt $(4 \mid 1)$, abgelesen aus dem eigenen Bild. Liegen die eigenen Geraden falsch, zählt der richtig abgelesene Schnittpunkt der eigenen Geraden.\\
(b) Bestätigen & 1 & Probe in beiden Gleichungen ($4 + 2 = 6$ \checkmark, $4 - 1 = 3$ \checkmark) oder ein Verfahren, zum Beispiel (I) $-$ (II): $3y = 3$, $y = 1$, $x = 4$. Folgezeile: Eine richtig gerechnete Bestätigung, die zeigt, dass der eigene abgelesene Punkt nicht stimmt, und die richtige Lösung $(4 \mid 1)$ nennt, gibt den Punkt.\\
\end{raster}
\fehler{Vorzeichen beim Auflösen, (I) als $y = 0.5x + 3$: Zeile (a) 0; die eigenen Geraden schneiden sich erst bei $(12 \mid 9)$, ausserhalb des Bildes — Ablesezeile 0; eine Rechnung, die $(4 \mid 1)$ ergibt, gibt die Bestätigungszeile $\to$ 1 von 3\,P.
Nur abgelesen, ohne Probe oder Rechnung: Bestätigungszeile 0.}
\end{lpaufgabe}

\begin{lpaufgabe}{G2}{Lösungsmenge und Lage}{3 P}
\begin{raster}
Eliminieren & 1 & $4 \cdot \text{(I)}$ ist genau (II): $2x + 4y = 16$. Also $4 \cdot \text{(I)} - \text{(II)}$: $0 = 0$, immer wahr. Gleichwertig: (II) ist das Vierfache von (I), so festgestellt.\\
Lösungsmenge & 1\E & $\mathbb{L} = \{(x \mid 4 - 0.5x) : x \in \mathbb{R}\}$. Gleichwertig: $\{(x \mid -\tfrac{1}{2}x + 4) : x \in \mathbb{R}\}$, $\{(8 - 2y \mid y) : y \in \mathbb{R}\}$ oder in Worten «alle Paare $(x \mid y)$ mit $y = 4 - 0.5x$».\\
Lage & 1 & Die Geraden sind identisch (fallen zusammen) — unendlich viele Lösungen.\\
\end{raster}
\fehler{$\mathbb{L} = \mathbb{R}$ oder «alle Zahlen»: Lösungsmengenzeile 0 — Lösungen sind Paare, und nur die auf der Geraden.
$0 = 0$ richtig hergeleitet, aber als Widerspruch gelesen, «keine Lösung, parallel»: Eliminierzeile 1, Lösungsmengen- und Lagezeile 0 (derselbe Begriffsfehler, kein Folgepunkt) $\to$ 1 von 3\,P.
Nur ein einzelnes Paar wie $(0 \mid 4)$ angegeben: Lösungsmengenzeile 0.}
\end{lpaufgabe}

\begin{lpaufgabe}{G3}{Lösungsfall am Bild}{2 P}
\begin{raster}
(a) Zuordnung & 1\E & (1) $0 = 5$ — Bild A (parallel, keine Lösung); (2) $0 = 0$ — Bild C (identisch, unendlich viele); (3) $x = 2$ — Bild B (die Geraden schneiden sich, bei $x = 2$). Alle drei nötig.\\
(b) Zwei Lösungen? & 1 & Nein, mit Begründung: Zwei Geraden haben keinen, genau einen oder unendlich viele gemeinsame Punkte. Gleichwertig: Haben zwei Geraden zwei gemeinsame Punkte, sind sie identisch, denn durch zwei Punkte geht nur eine Gerade — dann gibt es unendlich viele Lösungen.\\
\end{raster}
\fehler{Nur «Nein» ohne Begründung: Zeile (b) 0. Eine vertauschte Zuordnung: Zeile (a) 0.}
\end{lpaufgabe}

\begin{lpaufgabe}{G4}{Verfahren wählen}{3 P}
\begin{raster}
Verfahren & 1 & Gleichsetzen: Beide Gleichungen sind nach $x$ aufgelöst. Gleichwertig: Einsetzen ($x = 4 - y$ in die zweite Gleichung gibt dieselbe Gleichung), mit diesem Grund. Verfahren \emph{und} Grund nötig.\\
Gleichsetzen & 1\E & $4 - y = 2y - 5$, also $9 = 3y$, $y = 3$.\\
Lösung & 1 & Folgezeile: $x = 4 - 3 = 1$, $\mathbb{L} = \{(1 \mid 3)\}$.\\
\end{raster}
\fehler{Vorzeichen beim Hinüberbringen, $3y = -1$: zweite Zeile 0; weiter gerechnet $x = \tfrac{13}{3}$: Lösungszeile 1 (Folgezeile) $\to$ 2 von 3\,P.
Additionsverfahren gewählt (erst ordnen nötig): Verfahrenszeile 0; richtig gelöst: die beiden anderen Zeilen voll.
Koordinaten vertauscht, $\mathbb{L} = \{(3 \mid 1)\}$: Lösungszeile 0.}
\end{lpaufgabe}

\begin{lpaufgabe}{G5}{Erst ordnen}{3 P}
\begin{raster}
Ordnen & 1 & (I) $5x + 3y = 1$, (II) $2x - 5y = 19$. Beide nötig.\\
Vervielfachen und lösen & 1\E & Zum Beispiel $5 \cdot \text{(I)} + 3 \cdot \text{(II)}$: $25x + 15y + 6x - 15y = 5 + 57$, also $31x = 62$, $x = 2$ — beide Seiten vervielfacht. Gleichwertig: $x$ eliminieren mit $2 \cdot \text{(I)} - 5 \cdot \text{(II)}$: $31y = -93$, $y = -3$.\\
Zweiter Wert und Probe & 1 & Folgezeile: aus (I) $10 + 3y = 1$, $y = -3$, $\mathbb{L} = \{(2 \mid -3)\}$, \emph{und} die Probe in beiden Gleichungen: $3 \cdot (-3) = -9 = 1 - 10$ \checkmark, $2 \cdot 2 = 4 = 19 - 15$ \checkmark. Beides nötig.\\
\end{raster}
\fehler{Vorzeichen beim Ordnen, $-5x + 3y = 1$: Ordnungszeile 0; Lösung $\left(-\tfrac{62}{19} \mid -\tfrac{97}{19}\right)$: zweite Zeile 0; dritte Zeile nur mit einer Probe, die den Fehler aufdeckt $\to$ höchstens 1 von 3\,P.
Rechte Seite nicht vervielfacht ($25x + 15y = 1$): zweite Zeile 0; dritte Zeile nach der Proberegel $\to$ höchstens 2 von 3\,P.}
\end{lpaufgabe}

\begin{lpaufgabe}{G6}{Fehler finden}{2 P}
\begin{raster}
Fehler & 1 & Bei (I) $-$ (II) fällt $y$ nicht weg: $2y - (-2y) = 4y$, es entsteht $2x + 4y = 8$, nicht $2x = 8$. (Die Probe zeigt es: $(4 \mid 2)$ in (I) gibt $12 + 4 = 16 \neq 8$.) Eine Begründung ist nötig; gleichwertig ist die Probe allein, wenn sie zeigt, dass $(4 \mid 2)$ die Gleichung (I) nicht erfüllt.\\
Richtig gelöst & 1\E & (I) $+$ (II): $4x = 8$, $x = 2$; aus (II) $y = 1$. $\mathbb{L} = \{(2 \mid 1)\}$. Gleichwertig mit jedem anderen Verfahren.\\
\end{raster}
\fehler{Nur «falsch gerechnet» oder «es hätte addiert werden müssen» ohne Grund: Fehlerzeile 0.}
\end{lpaufgabe}

\begin{lpaufgabe}{G7}{Parameter in beiden Gleichungen}{3 P}
\begin{raster}
(a) gleiche Steigung & 1\E & (I) ist links ein Vielfaches von (II), wenn $\tfrac{k}{1} = \tfrac{4}{k}$, also $k^2 = 4$: $k = 2$ oder $k = -2$. Gleichwertig: Steigungen $-\tfrac{k}{4} = -\tfrac{1}{k}$; oder rechnerisch $k \cdot \text{(II)} - \text{(I)}$: $(k^2 - 4)\,y = 4k - 8$. Beide Werte nötig.\\
(b) Lösungsfälle & 1 & Folgezeile: $k = 2$: $2x + 4y = 8$ ist das Doppelte von $x + 2y = 4$, auch rechts — identische Geraden, unendlich viele Lösungen. $k = -2$: $-2x + 4y = 8$ ist links das $(-2)$-Fache von $x - 2y = 4$, rechts müsste $-8$ stehen — parallel, keine Lösung. Beide Fälle mit Begründung nötig. Wer in (a) nur einen oder einen falschen Wert hat und ihn richtig untersucht, erhält die Zeile.\\
(c) $k = 0$ & 1\E & $4y = 8$, $y = 2$, und $x = 4$: $\mathbb{L} = \{(4 \mid 2)\}$.\\
\end{raster}
\fehler{Nur $k = 2$ gefunden ($\sqrt4$ ohne $-2$): Zeile (a) 0; (b) für $k = 2$ allein: 0 (die Zeile verlangt beide Fälle, und der Fall $k = -2$ fehlt aus demselben Grund) $\to$ höchstens 1 von 3\,P.
Bei $k = 2$ «parallel» oder bei $k = -2$ «identisch»: Zeile (b) 0.}
\end{lpaufgabe}

\begin{lpaufgabe}{G8}{Drei Variablen}{6 P}
\begin{raster}
(a) Begründung & 1 & $z$ eliminieren, weil (III) schon ohne $z$ ist — eine einzige Elimination genügt. Gleichwertig: $y$ eliminieren (wie im Kontrollclip) oder $x$ aus (III) freistellen ($x = -5 - 3y$) und einsetzen, je mit Begründung.\\
(a) Elimination und 2×2 & 1\E & (II) $-$ (I): $x - 2y = 5$ (IV); mit (III): (III) $-$ (IV) gibt $5y = -10$, $y = -2$, $x = 1$. Gleichwertig auf einem anderen Weg.\\
(a) dritte Variable & 1 & Folgezeile: in (I) $1 - 2 + z = 3$, $z = 4$; $\mathbb{L} = \{(1 \mid -2 \mid 4)\}$.\\
(a) Probe & 1 & In allen drei Gleichungen: $1 - 2 + 4 = 3$ \checkmark, $2 + 2 + 4 = 8$ \checkmark, $1 - 6 = -5$ \checkmark. Nur zwei geprüft: 0.\\
(b) keine Lösung & 1\E & (I) $+$ (II) gibt $3x + 2z = 11$. Mit $3x + 2z = c$ folgt $0 = c - 11$: Für jedes $c \neq 11$ ein Widerspruch, keine Lösung.\\
(b) $c = 11$ & 1 & Dann bleibt $0 = 0$, und nirgends entsteht ein Widerspruch: unendlich viele Lösungen (eine Variable frei). Die Lösungsmenge ist nicht verlangt.\\
\end{raster}
\fehler{Erst $y$ eliminieren ist ein gültiger Weg: (I) $+$ (II) gibt $3x + 2z = 11$; dann $y$ auch aus (III) eliminieren, zum Beispiel $3 \cdot \text{(I)} - \text{(III)}$: $2x + 3z = 14$ — volle Punkte. Nur wer mit $3x + 2z = 11$ und (III) direkt weiterrechnet, obwohl in (III) noch $y$ steht, kommt nicht zu einem 2×2-System: Eliminationszeile 0, die übrigen Zeilen nur auf einem richtigen Weg danach.
Vorzeichen bei (II) $-$ (I), etwa $x = 5$ (das $-2y$ vergessen): Eliminationszeile 0; mit eigenen Werten richtig weitergerechnet: Folgezeile 1, Probe deckt den Fehler auf $\to$ Probezeile 1.
(b) «$c = 11$: keine Lösung» oder «für $c \neq 11$ unendlich viele»: beide (b)-Zeilen 0.}
\end{lpaufgabe}


\section*{4 · Selbsteinschätzung}
\begin{tabularx}{\linewidth}{@{}l X@{}}\toprule
22 – 25 P & Die geprüften Teile sitzen. Wo du Punkte verloren hast: das Kapitel dieser Aufgabe nochmals (Zuordnung unten).\\
17 – 21 P & Den schwächsten Teil nochmals: Tüfteln und Übungen des Kapitels, in dem du die meisten Punkte verloren hast.\\
11 – 16 P & Zurück zu den Kapiteln aller Aufgaben, in denen du Punkte verloren hast.\\
0 – 10 P & Zurück zu Kapitel 1 und von dort der Reihe nach weiter.\\\midrule
\multicolumn{2}{@{}p{\linewidth}@{}}{\small\textbf{Aufgabe $\to$ Kapitel:} G1 $\to$ Kapitel 1 (Lösung und Schnittpunkt), die Bestätigung Kapitel 2 oder 3; G2 $\to$ Kapitel 4 (Keine oder unendlich viele Lösungen); G3 $\to$ Kapitel 1 und 4; G4 $\to$ Kapitel 2 (Einsetzen und Gleichsetzen) und 3 (Verfahrenswahl); G5, G6 $\to$ Kapitel 3 (Additionsverfahren); G7 $\to$ Kapitel 4; G8 $\to$ Kapitel 5 (Drei Variablen), (b) auch Kapitel 4. \textbf{Kompetenzen:} K1 (lösen) G1(b), G2, G4–G8; K2 (grafisch veranschaulichen und deuten) G1, G2, G3, G7.}\\\bottomrule
\end{tabularx}
\end{document}
